\item Una onda sonora plana en aire a $20.0^\circ\text{C}$, con longitud de onda de $589\text{ mm}$, incide sobre una superficie uniforme de agua a $25.0^\circ\text{C}$ a un ángulo de incidencia de $13.0^\circ$. Determine:
\begin{enumerate}
    \item el ángulo de refracción para la onda sonora y
    \item la longitud de onda del sonido en el agua.
    \item Ahora considere un rayo de luz amarilla de sodio, con longitud de onda de $589\text{ nm}$ en el vacío, que incide desde el aire sobre una superficie uniforme de agua a un ángulo de incidencia de $13.0^\circ$. Determine el ángulo de refracción y
    \item la longitud de onda de la luz en el agua.
    \item Compare y contraste el comportamiento del sonido y de la luz en este problema.
\end{enumerate}